\subsection{Anchoring and invariant moments}

Fix one line as an anchor and rephase the other representatives to make all its incident Gram entries $+\mu$. Index the eight nonanchor vertices by $0,\ldots,7$. Write the signature matrix in \eqref{eq:signature} as
\[
 S=\begin{pmatrix}0&\one^T\\\one&T\end{pmatrix},\qquad
 T_{ii}=0,\quad T_{ij}=x_{ij}+i y_{ij}\quad(i<j),
\]
with $T=T^*$. The 56 real variables satisfy
\begin{equation}\label{eq:anchored}
 x_{ij}^2+y_{ij}^2=1,\qquad T\one=\beta\one,\qquad
 T^2=8I+\beta T-J,\qquad J=\one\one^T.
\end{equation}
The first equations give the diagonal of the last matrix equation. The remaining real equations consist of a circle, real and imaginary row sums, and real and imaginary off-diagonal tightness, together with relabelings of these five types. Conversely, the block equations recover \eqref{eq:signature} and hence \eqref{eq:gram}.

A truncated moment functional $L$ is linear on real polynomials of the given degree limit. The constraints are $L(1)=1$, $L(hm)=0$ for selected equation multiples, and $L(p^2)\ge0$ for all $p$ in a specified monomial span $V$. Here $h$ is an equation polynomial and $m$ is a monomial. The last condition is the positive semidefiniteness of the matrix with entries $L(uv)$ in a basis of $V$. This is the truncated moment construction of Lasserre \cite{lasserre}. The quartic model uses all monomials of degree at most two for $V$ and every equation multiple from \eqref{eq:anchored} of total degree at most four.

Relabeling the eight nonanchor vertices and conjugating all entries preserve \eqref{eq:anchored}. These operations act on monomials by signed permutations, with a sign for each reversal of a $y$-edge and for conjugation of each $y$ factor. We average $L$ over $S_8\times C_2$. A putative frame need only give point evaluations before averaging. Averaging takes place in moment space, irrespective of the frame's projective automorphism group. Invariance permits one coordinate per signed monomial orbit. An orbit whose representative can be mapped to its negative has moment zero. These rules also determine all unused entries of a stored moment matrix from the representative coordinates.

\subsection{Exact reduction of the positivity condition}

For $k=3,4,5,6$, let $V_k$ consist of all monomials of degree at most two and the cubic monomials involving at most $k$ nonanchor vertices. The vertex support of a monomial is the union of the endpoints of its factors, counting each endpoint only once. Thus $V_6$ contains every cubic, while the lower stages include specified subsets. The moment matrices for invariant $L$ commute with the orthogonal signed-permutation action in the monomial basis. This is the setting of symmetry reduction for semidefinite constraints \cite{gatermannparrilo}.

Let $s=\max(4,k)$ for $V_k$ and $s=4$ for the quartic basis. Take $H$ to be the permutation group on the last $8-s$ vertices. Positivity on $V$ is equivalent to positivity on $V^H$. To see this, suppose the full symmetric matrix has a negative eigenvalue and let $P$ be its spectral projection. There is a basis monomial $m$ with $Pm\ne0$. Relabel $m$ to use only the first $s$ vertices. The relabeled monomial is fixed by $H$, and its projection under $P$ remains nonzero and fixed by $H$ because $P$ commutes with the signed group action. This gives a negative direction in $V^H$. The reverse implication follows from restriction.

The transpositions $(0\ 1)$ and $(2\ 3)$ and the conjugation involution commute on $V^H$. Each has two possible eigenvalues, giving eight simultaneous character spaces. Invariance makes the mixed matrix entries between different characters zero. Within a character, a basis consists of signed orbit sums of monomials for the subgroup generated by $H$ and these three involutions. Orbits with inconsistent signs contribute no basis vector. The orbit sums are orthogonal in monomial coefficient space, and their squared norms are the orbit sizes. The moment blocks in these bases are integer linear combinations of the orbit moments. All eight blocks being positive semidefinite is equivalent to the full positivity condition on the chosen $V$.

\subsection{Sextic equations and structural kernels}

The sextic stages retain the quartic equations. They add $L(hm)=0$ for the five types of base equations in \eqref{eq:anchored}, with total degree five or six and multipliers on at most the first $\min(k,4)$ vertices, using representatives for $h$ whose designated row or edge is in this set. Invariance determines their relabelings. The internal summation indices of $h$ range over the full set of eight vertices. We also impose rank equations. All $5\times5$ minors of $32G$ vanish, as do all $4\times4$ minors of $32(G_{\mathrm{tail}}-\mu^2 J)$, where $G_{\mathrm{tail}}$ is the principal Gram submatrix complementary to the anchor. Indeed, the second matrix is a scalar multiple of the Schur complement of $G$ at its anchor, so it has rank three. We take real and imaginary parts, with all multipliers of degree at most one for the Gram minors and at most two for the Schur minors, modulo relabeling. These constraints are necessary even for nonsymmetric frames. The model uses both principal and nonprincipal minors.

For the three- and four-vertex stages we take products $hm\in V_k$, with $h$ ranging over relabelings of the base equations and $m$ over monomials. We project $hm$ into each tail-fixed character subspace and require its moment pairing with every polynomial in that character block to vanish. In the five-vertex stage we add the raw pairings $L((hm)v)=0$ for such products and monomials $v\in V_5$ to the inherited four-vertex equations. Invariance and linearity imply the projected pairings. All such pairings have degree at most six, and the projected directions are finite linear combinations of relabeled equations times monomials. Their point evaluations vanish at each solution of \eqref{eq:anchored}. The independent verifier uses the archived row provenance to check the imposed systems at each stage.

In a block basis $u_j=\sum_{n\in O_j}\epsilon_n n$, the projection of $hm$ has coefficient
\[
 q_j=\frac{\sum_{n\in O_j}\epsilon_n[n](hm)}{|O_j|}
\]
on $u_j$. Thus the linear cross equations $M_b q=0$ place each such $q$ in the kernel of the block matrix $M_b$. This gives a facial reduction \cite{borweinwolkowicz}. Choose a basis of these directions over $\Q(\beta)$ and a pivot set $P_b$ on which the basis has invertible coordinate restriction. The complementary coordinate vectors then span a direct complement. Every quadratic form value on the original block equals its value on the complementary component since the kernel has been imposed. Consequently the reduced block is exactly the principal submatrix on the nonpivot coordinates. In the three- and four-vertex stages we also archive orthogonal complements, computed using the orbit-size weights.

\subsection{Feasible witnesses and matrix sizes}

The exact witness coordinates belong to $\Q(\beta)$. These witnesses satisfy normalization and all the specified linear equations. The original quartic character blocks are checked by exact PSD elimination. We also verify 48 additional quartic identities obtained by changing the anchor where the products on both sides remain within the degree limit. The equality rank including normalization remains 99, and the exact ranks of the eight original quartic blocks are $24,14,10,14,10,14,10,10$. For the three-vertex stage the independent replay checks the original eight blocks and the 763 homogeneous baseline equations. The structural cross equations are checked separately on the same moment vector. For the four- and five-vertex witnesses, exact tests verify the inherited equations and positivity via the block congruences and principal restriction above. Appendix~\ref{app:replay} gives the artifact paths and commands.

\begin{table}[tb]
\centering
\small
\begin{tabular}{@{}lrrl@{}}
\toprule
Relaxation & Cubic monomials & Moment coordinates & Block orders after structural reduction\\
\midrule
Quartic & 0 & 107 & $72,45,34,38,34,38,25,24$\\
Three-vertex sextic & 2{,}576 & 2{,}547 & $74,55,46,50,46,50,35,39$\\
Four-vertex sextic & 14{,}056 & 2{,}716 & $128,99,86,98,86,98,77,75$\\
Five-vertex sextic & 27{,}496 & 2{,}719 & $373,342,284,302,284,302,239,235$\\
Six-vertex sextic & 30{,}856 & 2{,}719 & $886,855,702,721,702,721,579,573$\\
\bottomrule
\end{tabular}
\caption{Sizes of the five relaxations with exact feasible witnesses. The quartic row lists the eight original character blocks. The five-vertex row precedes its second facial reduction.}\label{tab:sizes}
\end{table}

For the four- and five-vertex stages, the change of coordinates $L(m)=\beta^{\deg(m)\bmod2}z_m$ makes the equality system rational after common factors are removed. The affine charts are constructed by rational elimination and checked against all equations. The rank evidence consists of invertible minors and exact identities for the free columns. The four-vertex equality rank including normalization is $2{,}346$, leaving $370$ free coordinates. The linear map from these directions into the eight PSD blocks has rank $97$. Since the remaining directions can be removed without changing the image, the numerical search can operate in 97 dimensions. The five-vertex stage adds 7{,}027 distinct cross equations, has equality rank $2{,}642$ including normalization, and leaves $77$ free coordinates with PSD image rank $75$. Here the extra equations can reduce the image dimension even though the matrix blocks are larger. Exact checks of rank use chart relations and finite-prime minors certified against denominator-cleared rational data. The floating-point barrier iterates are used to locate candidate witnesses. Exact equality and positivity checks establish the assertions in Theorem~\ref{thm:moment}.

\subsection{The common kernel of the five-vertex pencil}

Let $M_b(y)$ be the eight five-vertex blocks on the first reduced face, with $y$ the 77-dimensional affine chart coordinate. In the archive's character order their common kernels have dimensions $(9,6,4,6,4,6,3,3)$. They are verified by multiplication against the constant block and all 77 coefficient blocks. Restricting to complementary principal coordinates gives orders $(364,336,280,296,280,296,236,232)$. On each of these blocks the same exact feasible witness is positive definite. Consequently each common kernel is the full kernel at that witness.

Define the common margin as the supremum of $t$ for which some $y$ satisfies $M_b(y)-tI\succeq0$ for every $b$. Every $y$ has the archived null directions, so $t\le0$. The exact witness attains $t=0$. Positive diagonal rescaling of the blocks preserves this zero optimal margin, since it preserves the existence of common null directions and the PSD witness. After the second facial reduction, strict feasibility on the eight complementary blocks follows from their positive definiteness at the same witness.

\subsection{Six-vertex extension probe}\label{sec:six}

At six vertices the PSD test space contains all $30{,}856$ cubic monomials in the 56 variables. The probe retains the earlier choice of high-degree equality multipliers and adds the six-vertex structural cross equations. Its equality family is the one selected for the probe, with all cubics admitted in the PSD test space. It generates $99{,}821$ cross-product proposals and $7{,}048$ distinct nonzero rows. With the same $2{,}719$ moment coordinates, two additional independent equations increase the equality rank to $2{,}644$ including normalization. The resulting $75$ free directions have PSD image rank $75$. After structural facial reduction the eight block orders are $(886,855,702,721,702,721,579,573)$, with $2{,}105{,}310$ packed rows stored as sparse maps. The estimated peak storage for the blockwise Newton computation is $1{,}283{,}845{,}760$ bytes.

The exact five-vertex witness fails six of the new rows. For example, in the archive's zero-based parity-scaled coordinates, one of the new constraints sets $z_{90}-z_{2617}-z_{2618}$ to zero. At the five-vertex witness this expression equals $-242106776781/12500000000000$. The check fixes all $2{,}719$ inherited coordinates. Thus any six-vertex witness must move at least one coordinate of the archived five-vertex point.

The guarded computation used a 3{,}520 MiB memory limit and a total wall allowance of 1{,}790 seconds. The model assembly, elimination, and numerical search consumed 1{,}186.7 seconds of guarded wall time, with a peak sampled resident set of 768.8 MiB. After 35 Newton steps, the smallest eigenvalue among the eight unshifted, diagonally scaled blocks was approximately $3.06\times 10^{-8}$. The barrier shift at this iterate was about $-0.02095$. The unshifted blocks determine the positivity test for the returned moment vector.

We rounded the 75 free coordinates to rationals with denominator $10^{14}$ and evaluated the exact affine chart. The independent verifier checked the composition with the five-vertex chart, the new equations, and all $2{,}105{,}310$ packed stencils against the reconstructed symmetry bases. Conjugation by the diagonal degree-parity scaling turns each block of the witness into a rational matrix. An integer congruence gives strict diagonal dominance with positive diagonal in each full structurally reduced block. The exact witness therefore proves strict feasibility on the six-vertex structural face. A separate full-rank modular specialization verifies that the eight affine blocks have trivial common kernels. This completes Theorem~\ref{thm:moment}.
